\documentclass[12pt]{article}
\usepackage[utf8]{inputenc}
\usepackage[T1]{fontenc}
\usepackage[margin=1in]{geometry}
\usepackage{setspace}
\usepackage{xcolor}
\usepackage[colorlinks=true, urlcolor=blue!60!black, linkcolor=blue!60!black]{hyperref}
\setlength{\parindent}{2em}
\setlength{\parskip}{1.25em}
\renewcommand{\baselinestretch}{1.0}
\title{\textit{Trump v. Mazars USA, LLP}, 591 U.S. 848 (2020)}
\author{Case Brief}
\date{}

\begin{document}
\maketitle

\section{Case Metadata}
\begin{description}
  \item[Citation:] 591 U.S. 848 (2020) (Docket No. 19-715; consolidated with No. 19-760, \textit{Trump v. Deutsche Bank AG}).
  \item[Decided:] July 9, 2020 (Argued May 12, 2020).
  \item[Vote:] 7--2 (vacated and remanded).
  \item[Author:] Chief Justice John G. Roberts, Jr. (joined by Justices Ginsburg, Breyer, Sotomayor, Kagan, Gorsuch, and Kavanaugh).
  \item[Concurrences / Dissents:] None concurring. Justice Thomas filed a dissenting opinion; Justice Alito filed a dissenting opinion.
  \item[Chapter / Topic:] Separation of Powers and the Executive --- Congressional Investigatory Power and Presidential Records.
  \item[Primary Source URLs:] \href{https://supreme.justia.com/cases/federal/us/591/19-715/}{Justia (591 U.S. 848)}; \href{https://www.law.cornell.edu/supremecourt/text/19-715}{Cornell LII (No. 19-715)}; \href{https://www.supremecourt.gov/opinions/19pdf/19-715_febh.pdf}{SupremeCourt.gov Slip Opinion}.
\end{description}

\section{Facts and Procedural History}
Over five days in April 2019, three committees of the U.S. House of Representatives issued four subpoenas seeking roughly a decade of financial records concerning President Donald J. Trump, his children, and affiliated business entities:
\begin{itemize}
  \item \textbf{House Committee on Financial Services (April 11, 2019):} Issued two subpoenas---one to Deutsche Bank (seeking account activity, due diligence, foreign transactions, statements of net worth, tax returns, and suspicious activity from 2010 onward, with no time limit on certain records) and one to Capital One (covering more than a dozen affiliated entities from 2016 onward)---citing H. Res. 206 and committee oversight plans aimed at anti-money-laundering, counterterrorism, and real-estate transparency legislation.
  \item \textbf{House Permanent Select Committee on Intelligence (April 11, 2019):} Issued an identical subpoena to Deutsche Bank as part of an inquiry into Russian interference in the 2016 election and potential foreign financial leverage over the President or his associates, aimed at national-security and election-integrity legislation.
  \item \textbf{House Committee on Oversight and Reform (April 15, 2019):} Subpoenaed the President's personal accounting firm, Mazars USA, LLP, for 2011--2018 financial statements, auditors' reports, underlying source documents, communications, and engagement contracts without regard to time. Following testimony by former personal attorney Michael Cohen, Chairman Elijah Cummings cited four investigative aims: potential illegal conduct, undisclosed conflicts of interest, compliance with the Emoluments Clauses, and accuracy of federal financial disclosures.
\end{itemize}

President Trump (suing in his personal capacity), his children, and affiliated businesses filed two suits seeking declaratory and injunctive relief: \textit{Mazars} (No. 19-715) in the U.S. District Court for the District of Columbia and \textit{Deutsche Bank} (No. 19-760) in the U.S. District Court for the Southern District of New York. Petitioners did not assert executive privilege; instead, they argued that the subpoenas lacked a legitimate legislative purpose and violated the separation of powers. The third-party custodians took no position, and the House committees intervened.

Both challenges failed below. In \textit{Mazars}, the D.C. District Court granted judgment for the House (380 F. Supp. 3d 76), and a divided D.C. Circuit affirmed (940 F.3d 710), holding the Oversight subpoena served a valid legislative purpose relating to presidential financial disclosure laws; Judge Rao dissented, arguing that investigating presidential wrongdoing belongs exclusively to the impeachment power. In \textit{Deutsche Bank}, the S.D.N.Y. denied a preliminary injunction, and the Second Circuit affirmed in substantial part (943 F.3d 627), concluding that separation-of-powers concerns were not implicated because the records were personal and that the House could appropriately use the President as a ``case study'' for banking and national-security legislation; Judge Livingston dissented in part. The Supreme Court granted certiorari in both cases, consolidated them, and stayed the judgments below.

\section{Issue(s) Presented}
Whether congressional committees exceed their constitutional authority under Article I and the separation of powers when they issue legislative subpoenas to third-party custodians for a sitting President's nonprivileged personal financial records, and what legal standard governs judicial review of such subpoenas.

\section{Holding}
Vacated and remanded (7--2). The lower courts failed to take adequate account of the significant separation-of-powers concerns implicated by congressional subpoenas for a sitting President's personal information. Rejecting both the heightened executive-privilege standard urged by the President and the deferential legislative-purpose standard urged by the House, the Court held that courts must apply a balanced, four-factor (non-exhaustive) separation-of-powers inquiry before enforcing a congressional subpoena for presidential records.

\section{Reasoning}
\textbf{1. Historical Practice of Interbranch Accommodation.} Chief Justice Roberts emphasized that for more than two centuries---from President Washington's handling of the 1792 St. Clair inquiry and President Jefferson's response to the 1807 Burr conspiracy request through the Reagan and Clinton administrations---Congress and the President had resolved document disputes through the ``hurly-burly, the give-and-take of the political process'' without Supreme Court intervention. Under \textit{NLRB v. Noel Canning}, 573 U.S. 513 (2014), this longstanding practice is a consideration of great weight, imposing a duty of care not to disturb established interbranch working arrangements.

\textbf{2. Scope and Limits of Congressional Inquiry.} Although Article I contains no express investigative clause, each House possesses an implied power ``to secure needed information'' as an auxiliary to legislating (\textit{McGrain v. Daugherty}, 273 U.S. 135 (1927); \textit{Watkins v. United States}, 354 U.S. 178 (1957)). That power is bounded: a subpoena must advance a ``valid legislative purpose'' concerning a subject on which legislation could be had, may not be used for ``law enforcement'' (assigned to the Executive and Judiciary) or ``to expose for the sake of exposure,'' and remains subject to constitutional rights and common-law or executive privileges.

\textbf{3. Rejection of the President's and Solicitor General's Standard.} The President and Solicitor General argued that the House should satisfy the demanding standards of \textit{United States v. Nixon}, 418 U.S. 683 (1974) (``demonstrated, specific need'') and \textit{Senate Select Committee on Presidential Campaign Activities v. Nixon}, 498 F.2d 725 (D.C. Cir. 1974) (``demonstrably critical''). The Court refused to ``transplant that protection root and branch'' from privileged Oval Office communications to nonprivileged, private records. Applying executive-privilege thresholds to all presidential information would seriously impede Congress's constitutional duty to legislate effectively.

\textbf{4. Rejection of the House's Deferential Standard.} Equally, the House's reliance on ordinary legislative-subpoena precedents (\textit{Barenblatt}, \textit{Eastland}) ignored the reality that Congress and the President are ``opposite and rival'' political branches in an ongoing institutional relationship. Because Congress's legislative jurisdiction is vast, virtually any personal paper---financial, medical, or educational---could ``relate to'' conceivable legislation, giving Congress a limitless subpoena power to harass or render the President ``complaisan[t] to the humors of the Legislature'' and allowing it to walk away from the political bargaining table. Nor do separation-of-powers concerns vanish because the papers are personal (the President alone composes a branch of government, linking ``the interest of the man'' with ``the constitutional rights of the place'') or because the subpoenas target third-party custodians (the Constitution ``deals with substance, not shadows'').

\textbf{5. The Four-Factor \textit{Mazars} Balancing Framework.} Steering a middle course, the Court directed lower courts evaluating congressional subpoenas for presidential information to weigh at least four special considerations:
\begin{enumerate}
  \item \textbf{Necessity / Alternative Sources:} Carefully assess whether the asserted legislative purpose warrants the significant step of involving the President and his papers. Congress may not rely on the President's information if other sources could reasonably supply what it needs, and it may not treat the President as a ``case study'' for general legislation.
  \item \textbf{Tailoring / Breadth:} Insist on a subpoena no broader than reasonably necessary to support Congress's legislative objective, safeguarding against unnecessary intrusion into the Office of the President.
  \item \textbf{Evidence of Legislative Purpose:} Scrutinize the nature of the evidence offered by Congress; the more detailed and substantial the evidence of a valid legislative purpose, the better---especially when Congress contemplates constitutionally sensitive legislation regulating the Presidency itself.
  \item \textbf{Burdens on the President:} Assess the institutional burdens imposed on the President, recognizing that congressional subpoenas come from a rival political branch with ongoing incentives to seek institutional advantage.
\end{enumerate}

\section{Concurring and Dissenting Opinions}
\begin{itemize}
  \item \textbf{Justice Thomas (Dissenting):} Argued from original meaning and founding-era practice that Congress possesses \textit{no} implied legislative power to subpoena private, nonofficial documents---whether belonging to the President or any private citizen. Distinguishing Parliament's supremacy from Congress's enumerated powers and criticizing \textit{McGrain v. Daugherty} as unmoored from text and history (preferring the skepticism of \textit{Kilbourn v. Thompson}, 103 U.S. 168 (1881)), Thomas concluded that if the House wishes to compel private presidential papers to investigate wrongdoing, it must invoke its enumerated impeachment power under Art. I, \S~2, cl. 5.
  \item \textbf{Justice Alito (Dissenting):} Assumed arguendo that legislative subpoenas for a President's personal records are not categorically barred, but viewed them as ``inherently suspicious'' and found ``disturbing evidence of an improper law enforcement purpose'' here. While agreeing that the lower courts erred, Alito deemed the majority's four-factor remand terms inadequate and would have barred enforcement unless the House specifically described the contemplated legislation, spelled out its constitutional authority to enact it, justified the subpoenas' scope, and proved why alternative sources were insufficient.
\end{itemize}

\section{Key Quotations}
\begin{quote}
``Congress and the President---the two political branches established by the Constitution---have an ongoing relationship that the Framers intended to feature both rivalry and reciprocity.'' --- Chief Justice Roberts, \textit{Trump v. Mazars USA, LLP}, 591 U.S. 848 (2020) (\href{https://supreme.justia.com/cases/federal/us/591/19-715/}{Justia}).
\end{quote}
\begin{quote}
``The President is the only person who alone composes a branch of government. As a result, there is not always a clear line between his personal and official affairs.'' --- Chief Justice Roberts (\href{https://supreme.justia.com/cases/federal/us/591/19-715/}{Justia}).
\end{quote}
\begin{quote}
``The President's unique constitutional position means that Congress may not look to him as a `case study' for general legislation.'' --- Chief Justice Roberts (\href{https://supreme.justia.com/cases/federal/us/591/19-715/}{Justia}).
\end{quote}
\begin{quote}
``Other considerations may be pertinent as well; one case every two centuries does not afford enough experience for an exhaustive list.'' --- Chief Justice Roberts (\href{https://supreme.justia.com/cases/federal/us/591/19-715/}{Justia}).
\end{quote}
\begin{quote}
``I would hold that Congress has no power to issue a legislative subpoena for private, nonofficial documents---whether they belong to the President or not. Congress may be able to obtain these documents as part of an investigation of the President, but to do so, it must proceed under the impeachment power.'' --- Justice Thomas, dissenting (\href{https://supreme.justia.com/cases/federal/us/591/19-715/}{Justia}).
\end{quote}

\section{Doctrinal Significance and Modern Legacy (Through 2026)}
\textit{Mazars} is the foundational Supreme Court decision governing congressional subpoenas for presidential records and a centerpiece of the Roberts Court's structural separation-of-powers jurisprudence. Decided the same day as \href{https://www.law.cornell.edu/supremecourt/text/19-635}{\textit{Trump v. Vance}, 591 U.S. 786 (2020)} (which rejected absolute immunity or a heightened-need standard for a state criminal grand jury subpoena directed to Mazars for the same records), \textit{Mazars} draws a sharp institutional line between judicial/criminal process and interbranch legislative oversight: courts treat congressional demands on the President with far greater structural skepticism because Congress is a rival political branch.

On remand in \href{https://law.justia.com/cases/federal/appellate-courts/cadc/21-5176/21-5176-2022-07-08.html}{\textit{Trump v. Mazars USA, LLP}, 39 F.4th 774 (D.C. Cir. 2022)}, the D.C. Circuit held that heightened separation-of-powers concerns persisted even after Donald Trump left office, applied the four \textit{Mazars} factors, upheld the Oversight Committee's authority to obtain certain 2016--2018 financial records in aid of presidential-ethics legislation, and narrowed the subpoena's scope; the parties subsequently settled on August 30, 2022, leading to document production and dismissal with prejudice in July 2023 (\href{https://clearinghouse.net/case/17222/}{Civil Rights Litigation Clearinghouse}). Practically, however, the multi-year litigation demonstrated that \textit{Mazars}'s open-ended balancing test enables a President to delay congressional enforcement well past the end of a two-year Congress.

Through 2026, \textit{Mazars} sits within a broader, two-sided separation-of-powers architecture constructed by the Roberts Court. On the executive-protection side, \href{https://supreme.justia.com/cases/federal/us/603/23-939/}{\textit{Trump v. United States}, 603 U.S. 593 (2024)}, explicitly echoed \textit{Mazars}'s caution that ``[o]ne case'' in more than ``two centuries does not afford enough experience'' when recognizing absolute presidential criminal immunity for core constitutional acts and at least presumptive immunity for other official acts; and \href{https://supreme.justia.com/cases/federal/us/609/25-332/}{\textit{Trump v. Slaughter}, 609 U.S. \_\_\_ (2026)}, reinforced the President's removal authority under Article II. Conversely, on the executive-constraint side, the Court has insisted that the Executive cannot bypass Article I lawmaking or taxing prerogatives, invoking the major questions doctrine and strict statutory limits to invalidate executive overreach in \href{https://supreme.justia.com/cases/federal/us/597/20-1530/}{\textit{West Virginia v. EPA}, 597 U.S. 697 (2022)}, \href{https://supreme.justia.com/cases/federal/us/600/22-506/}{\textit{Biden v. Nebraska}, 600 U.S. 477 (2023)}, and \href{https://supreme.justia.com/cases/federal/us/607/24-1287/}{\textit{Learning Resources, Inc. v. Trump}, 607 U.S. 229 (2026)} (striking down unilateral presidential tariffs imposed under IEEPA). Together, these decisions reflect a Court actively policing formal interbranch boundaries---checking congressional encroachment on the Presidency in \textit{Mazars} while checking presidential encroachment on Congress's legislative and taxing powers in \textit{Nebraska} and \textit{Learning Resources}.

\section{Pedagogical Critiques and Exam Hypotheticals}
\textbf{Pedagogical Critiques:}
\begin{itemize}
  \item \textbf{Indeterminacy and Delay as Functional Victory:} Critics note that the Court's non-exhaustive, four-factor balancing test offers little bright-line guidance, inviting protracted district-court and appellate litigation that routinely outlasts the two-year lifespan of a House subpoena.
  \item \textbf{Asymmetry with \textit{Trump v. Vance}:} Students should interrogate why a single locally elected county prosecutor in \textit{Vance} faced a lower constitutional bar to subpoena the exact same Mazars records than three committees of the U.S. House of Representatives in \textit{Mazars}.
  \item \textbf{The Impeachment-Channeling Problem:} Justice Thomas's dissent forces a structural question: if Congress can compel private presidential papers only via impeachment, does that incentivize premature, hyper-partisan impeachment inquiries whenever Congress needs facts to evaluate ethics or disclosure statutes?
\end{itemize}

\textbf{Exam Hypothetical:}
Suppose the House Committee on Ways and Means is considering a bill to cap foreign-government hotel and resort payments to any federal official, including the President. The Committee issues a subpoena to a private hospitality management firm for all booking ledgers, guest lists, and pricing contracts from 2018--2026 for a resort owned by the sitting President's family trust, asserting that ``examining the President's resort operations will illuminate industry loopholes.'' The President sues in his personal capacity to enjoin the firm from complying, without invoking executive privilege.
\begin{itemize}
  \item \textbf{Majority / \textit{Mazars} Analysis:} The President's personal capacity and the third-party custodian do not eliminate separation-of-powers scrutiny. Under Factor 1, the House cannot use the President as a ``case study'' for general hospitality or ethics legislation if public filings, other officials' disclosures, or industry data could reasonably supply the needed information. Under Factor 2, an eight-year demand for all guest lists and contracts is likely overbroad relative to foreign-government payments. Under Factor 3, because regulating the President raises sensitive constitutional questions (e.g., Emoluments Clause enforcement), the Committee must offer detailed, substantial evidence tying these specific records to a valid legislative aim rather than disguised law enforcement. Under Factor 4, the court must weigh the institutional burden and partisan-leverage risk.
  \item \textbf{``Opposite Side'' Contrast:} Compare this scenario to a \textit{different} subpoena in which (a) the House formally opens an impeachment inquiry into alleged bribery involving foreign bookings at the resort (triggering the impeachment channel that even Justice Thomas recognized in \textit{Mazars}), or (b) a federal or state grand jury subpoenas the exact same resort records in a criminal tax-fraud investigation (governed by \textit{Trump v. Vance}, where \textit{Mazars}'s four-factor interbranch test does not apply and the President receives only ordinary citizen protections plus Supremacy Clause / Article II impairment defenses). Contrasting the three channels---legislative oversight (\textit{Mazars}), impeachment (Thomas dissent), and criminal process (\textit{Vance})---tests whether students understand that the constitutional standard turns on the \textit{institutional actor and constitutional function}, not merely the privacy of the underlying documents.
\end{itemize}

\end{document}
